% --- Archon physics formalization source begin ---
% archon:physics
% archon:covers PhyXMiniProblems/problem_phyx_mini_0451.lean
% archon:source-report reports/phyx_mini/problem_phyx_mini_0451.source.json

\chapter{Physics problem phyx\_mini\_0451}
\label{ch:PhyXMiniProblems_problem_phyx_mini_0451}

\paragraph{Problem source.}
\#\# Physical scenario

Two kilograms of water at \$120 \textbackslash{}, \textasciicircum{}\{\textbackslash{}circ\}\textbackslash{}text\{C\}\$ with a quality of \$25\textbackslash{}\%\$ has its temperature raised \$20 \textbackslash{}, \textasciicircum{}\{\textbackslash{}circ\}\textbackslash{}text\{C\}\$ in a constant-volume process as in figure.

\#\# Question

What is the heat transfer in the process?

\#\# Answer choices

- A: A: 1224.38 kJ

- B: B: 1009.92 kJ

- C: C: 503.48 kJ

- D: D: 877.8 kJ

\#\# Figure caption (auxiliary; use the image as primary evidence)

The image appears to be a diagram illustrating a heating mechanism, likely for a boiler or a similar device. Here's a detailed description:

1. **Container/Boiler**: A rectangular container is depicted. Inside the container, there is a body of liquid, represented by horizontal blue lines at the bottom.

2. **Heat Source**: Beneath the container, there are four flames, indicating a heat source. Each flame is centralized over a rectangular base, evenly spaced along the bottom.

3. **Pipe and Valve**: On top of the container, a pipe leads upward and then bends to the right. Just above where the pipe emerges from the container, there is a circular symbol with a cross inside, likely representing a valve controlling the flow through the pipe.

Overall, the setup suggests a typical heating and steam generation configuration, often used in boilers or similar systems.

\#\# Dataset metadata

Temperature and Heat Transfer; Physical Model Grounding Reasoning, Multi-Formula Reasoning

\paragraph{Recorded answer/context.}
D: D: 877.8 kJ

\paragraph{Figure/image path.}
/root/proposal\_for\_physic/science-mango/phyx\_mini\_run/phyx\_data/test\_image/451.png

\paragraph{Formalization target.}
create a compiling Lean file with sorry bodies at `PhyXMiniProblems/problem\_phyx\_mini\_0451.lean`. The Lean declarations must preserve the physical quantities, dimensions or dimensional roles, figure labels, governing-law hypotheses, and final relation expressed by this problem.
Use Mathlib/Physlib names found through LeanExplore where available. If a physics API is missing, introduce faithful local abstractions rather than scalar placeholder aliases.

\begin{theorem}[Physics formalization target]
\label{thm:physics:phyx_mini_0451:target}
\lean{PhyXMiniProblems.ProblemPhyXMini0451.heat_transfer_to_rigid_water_vessel}
\uses{def:physics:phyx-mini-0451:phyxminiproblems-problemphyxmini0451-energyinkilojoules, def:physics:phyx-mini-0451:phyxminiproblems-problemphyxmini0451-rigidwatervesselsetup, def:physics:phyx-mini-0451:phyxminiproblems-problemphyxmini0451-matchesrigidwatervesselscenario, def:physics:phyx-mini-0451:phyxminiproblems-problemphyxmini0451-matchessuppliedfigure, def:physics:phyx-mini-0451:phyxminiproblems-problemphyxmini0451-matchesproblemreadouts, def:physics:phyx-mini-0451:phyxminiproblems-problemphyxmini0451-matchesreferencesaturatedwaterdata, def:physics:phyx-mini-0451:phyxminiproblems-problemphyxmini0451-satisfiesrigidvesselwaterlaws, def:physics:phyx-mini-0451:phyxminiproblems-problemphyxmini0451-answerchoice, def:physics:phyx-mini-0451:phyxminiproblems-problemphyxmini0451-isreportedheatchoice}
The assigned autoformalize agent should translate this physics problem into Lean declarations in the covered file, with theorem and lemma proof bodies written as `by sorry`.
\end{theorem}
\begin{proof}
This is an autoformalization task, not a proof task. Produce statements that can later be proved by the physics prover without weakening the physical model.
\end{proof}

% --- Archon declaration topology begin ---
\section{Declaration topology}
\begin{definition}[volume Dimension]
  \label{def:physics:phyx-mini-0451:phyxminiproblems-problemphyxmini0451-volumedimension}
  \lean{PhyXMiniProblems.ProblemPhyXMini0451.volumeDimension}
  The physical dimension of volume, L³
\end{definition}
\begin{proof}
  This is the declared model definition of volume Dimension.
\end{proof}
\begin{definition}[specific Volume Dimension]
  \label{def:physics:phyx-mini-0451:phyxminiproblems-problemphyxmini0451-specificvolumedimension}
  \lean{PhyXMiniProblems.ProblemPhyXMini0451.specificVolumeDimension}
  \uses{def:physics:phyx-mini-0451:phyxminiproblems-problemphyxmini0451-volumedimension}
  The physical dimension of specific volume, L³ M⁻¹
\end{definition}
\begin{proof}
  This is the declared model definition of specific Volume Dimension.
\end{proof}
\begin{definition}[specific Internal Energy Dimension]
  \label{def:physics:phyx-mini-0451:phyxminiproblems-problemphyxmini0451-specificinternalenergydimension}
  \lean{PhyXMiniProblems.ProblemPhyXMini0451.specificInternalEnergyDimension}
  The physical dimension of specific internal energy, L² T⁻²
\end{definition}
\begin{proof}
  This is the declared model definition of specific Internal Energy Dimension.
\end{proof}
\begin{definition}[Mass Quantity]
  \label{def:physics:phyx-mini-0451:phyxminiproblems-problemphyxmini0451-massquantity}
  \lean{PhyXMiniProblems.ProblemPhyXMini0451.MassQuantity}
  A nonnegative, unit-independent physical mass
\end{definition}
\begin{proof}
  This is the declared model definition of Mass Quantity.
\end{proof}
\begin{definition}[Volume Quantity]
  \label{def:physics:phyx-mini-0451:phyxminiproblems-problemphyxmini0451-volumequantity}
  \lean{PhyXMiniProblems.ProblemPhyXMini0451.VolumeQuantity}
  \uses{def:physics:phyx-mini-0451:phyxminiproblems-problemphyxmini0451-volumedimension}
  A nonnegative, unit-independent physical volume
\end{definition}
\begin{proof}
  This is the declared model definition of Volume Quantity.
\end{proof}
\begin{definition}[Specific Volume Quantity]
  \label{def:physics:phyx-mini-0451:phyxminiproblems-problemphyxmini0451-specificvolumequantity}
  \lean{PhyXMiniProblems.ProblemPhyXMini0451.SpecificVolumeQuantity}
  \uses{def:physics:phyx-mini-0451:phyxminiproblems-problemphyxmini0451-specificvolumedimension}
  A nonnegative, unit-independent thermodynamic specific volume
\end{definition}
\begin{proof}
  This is the declared model definition of Specific Volume Quantity.
\end{proof}
\begin{definition}[Specific Internal Energy Quantity]
  \label{def:physics:phyx-mini-0451:phyxminiproblems-problemphyxmini0451-specificinternalenergyquantity}
  \lean{PhyXMiniProblems.ProblemPhyXMini0451.SpecificInternalEnergyQuantity}
  \uses{def:physics:phyx-mini-0451:phyxminiproblems-problemphyxmini0451-specificinternalenergydimension}
  A signed, unit-independent specific internal energy. Signed values permit an arbitrary thermodynamic reference datum; the tabulated states in this problem have positive readouts
\end{definition}
\begin{proof}
  This is the declared model definition of Specific Internal Energy Quantity.
\end{proof}
\begin{definition}[mass In Kilograms]
  \label{def:physics:phyx-mini-0451:phyxminiproblems-problemphyxmini0451-massinkilograms}
  \lean{PhyXMiniProblems.ProblemPhyXMini0451.massInKilograms}
  \uses{def:physics:phyx-mini-0451:phyxminiproblems-problemphyxmini0451-massquantity}
  Read a physical mass in kilograms
\end{definition}
\begin{proof}
  This is the declared model definition of mass In Kilograms.
\end{proof}
\begin{definition}[volume In Cubic Metres]
  \label{def:physics:phyx-mini-0451:phyxminiproblems-problemphyxmini0451-volumeincubicmetres}
  \lean{PhyXMiniProblems.ProblemPhyXMini0451.volumeInCubicMetres}
  \uses{def:physics:phyx-mini-0451:phyxminiproblems-problemphyxmini0451-volumequantity}
  Read a physical volume in cubic metres
\end{definition}
\begin{proof}
  This is the declared model definition of volume In Cubic Metres.
\end{proof}
\begin{definition}[specific Volume In Cubic Metres Per Kilogram]
  \label{def:physics:phyx-mini-0451:phyxminiproblems-problemphyxmini0451-specificvolumeincubicmetresperkilogram}
  \lean{PhyXMiniProblems.ProblemPhyXMini0451.specificVolumeInCubicMetresPerKilogram}
  \uses{def:physics:phyx-mini-0451:phyxminiproblems-problemphyxmini0451-specificvolumequantity}
  Read a specific volume in cubic metres per kilogram
\end{definition}
\begin{proof}
  This is the declared model definition of specific Volume In Cubic Metres Per Kilogram.
\end{proof}
\begin{definition}[specific Internal Energy In Joules Per Kilogram]
  \label{def:physics:phyx-mini-0451:phyxminiproblems-problemphyxmini0451-specificinternalenergyinjoulesperkilogram}
  \lean{PhyXMiniProblems.ProblemPhyXMini0451.specificInternalEnergyInJoulesPerKilogram}
  \uses{def:physics:phyx-mini-0451:phyxminiproblems-problemphyxmini0451-specificinternalenergyquantity}
  Read a specific internal energy in joules per kilogram
\end{definition}
\begin{proof}
  This is the declared model definition of specific Internal Energy In Joules Per Kilogram.
\end{proof}
\begin{definition}[specific Internal Energy In Kilojoules Per Kilogram]
  \label{def:physics:phyx-mini-0451:phyxminiproblems-problemphyxmini0451-specificinternalenergyinkilojoulesperkilogram}
  \lean{PhyXMiniProblems.ProblemPhyXMini0451.specificInternalEnergyInKilojoulesPerKilogram}
  \uses{def:physics:phyx-mini-0451:phyxminiproblems-problemphyxmini0451-specificinternalenergyquantity, def:physics:phyx-mini-0451:phyxminiproblems-problemphyxmini0451-specificinternalenergyinjoulesperkilogram}
  Read a specific internal energy in kilojoules per kilogram
\end{definition}
\begin{proof}
  This is the declared model definition of specific Internal Energy In Kilojoules Per Kilogram.
\end{proof}
\begin{definition}[energy In Joules]
  \label{def:physics:phyx-mini-0451:phyxminiproblems-problemphyxmini0451-energyinjoules}
  \lean{PhyXMiniProblems.ProblemPhyXMini0451.energyInJoules}
  Read heat, work, or total internal energy in joules
\end{definition}
\begin{proof}
  This is the declared model definition of energy In Joules.
\end{proof}
\begin{definition}[energy In Kilojoules]
  \label{def:physics:phyx-mini-0451:phyxminiproblems-problemphyxmini0451-energyinkilojoules}
  \lean{PhyXMiniProblems.ProblemPhyXMini0451.energyInKilojoules}
  \uses{def:physics:phyx-mini-0451:phyxminiproblems-problemphyxmini0451-energyinjoules}
  Read heat, work, or total internal energy in kilojoules
\end{definition}
\begin{proof}
  This is the declared model definition of energy In Kilojoules.
\end{proof}
\begin{definition}[Measured Temperature]
  \label{def:physics:phyx-mini-0451:phyxminiproblems-problemphyxmini0451-measuredtemperature}
  \lean{PhyXMiniProblems.ProblemPhyXMini0451.MeasuredTemperature}
  An absolute temperature together with the zero-preserving unit in which its magnitude is stored. The affine Celsius offset is introduced only in the scalar readout below
\end{definition}
\begin{proof}
  This is the declared model definition of Measured Temperature.
\end{proof}
\begin{definition}[temperature In Kelvins]
  \label{def:physics:phyx-mini-0451:phyxminiproblems-problemphyxmini0451-temperatureinkelvins}
  \lean{PhyXMiniProblems.ProblemPhyXMini0451.temperatureInKelvins}
  \uses{def:physics:phyx-mini-0451:phyxminiproblems-problemphyxmini0451-measuredtemperature}
  Kelvin readout of an absolute physical temperature
\end{definition}
\begin{proof}
  This is the declared model definition of temperature In Kelvins.
\end{proof}
\begin{definition}[temperature In Degrees Celsius]
  \label{def:physics:phyx-mini-0451:phyxminiproblems-problemphyxmini0451-temperatureindegreescelsius}
  \lean{PhyXMiniProblems.ProblemPhyXMini0451.temperatureInDegreesCelsius}
  \uses{def:physics:phyx-mini-0451:phyxminiproblems-problemphyxmini0451-measuredtemperature, def:physics:phyx-mini-0451:phyxminiproblems-problemphyxmini0451-temperatureinkelvins}
  Celsius readout, using the exact offset 273.15 K = 5463 / 20 K
\end{definition}
\begin{proof}
  This is the declared model definition of temperature In Degrees Celsius.
\end{proof}
\begin{definition}[Working Fluid]
  \label{def:physics:phyx-mini-0451:phyxminiproblems-problemphyxmini0451-workingfluid}
  \lean{PhyXMiniProblems.ProblemPhyXMini0451.WorkingFluid}
  The working substance in the depicted vessel
\end{definition}
\begin{proof}
  This is the declared model definition of Working Fluid.
\end{proof}
\begin{definition}[Process Stage]
  \label{def:physics:phyx-mini-0451:phyxminiproblems-problemphyxmini0451-processstage}
  \lean{PhyXMiniProblems.ProblemPhyXMini0451.ProcessStage}
  The two endpoint states of the heating process
\end{definition}
\begin{proof}
  This is the declared model definition of Process Stage.
\end{proof}
\begin{definition}[Water Phase Region]
  \label{def:physics:phyx-mini-0451:phyxminiproblems-problemphyxmini0451-waterphaseregion}
  \lean{PhyXMiniProblems.ProblemPhyXMini0451.WaterPhaseRegion}
  Thermodynamic phase region of an equilibrium water state
\end{definition}
\begin{proof}
  This is the declared model definition of Water Phase Region.
\end{proof}
\begin{definition}[Process Kind]
  \label{def:physics:phyx-mini-0451:phyxminiproblems-problemphyxmini0451-processkind}
  \lean{PhyXMiniProblems.ProblemPhyXMini0451.ProcessKind}
  Idealized type of the process specified in the prose
\end{definition}
\begin{proof}
  This is the declared model definition of Process Kind.
\end{proof}
\begin{definition}[Valve Status]
  \label{def:physics:phyx-mini-0451:phyxminiproblems-problemphyxmini0451-valvestatus}
  \lean{PhyXMiniProblems.ProblemPhyXMini0451.ValveStatus}
  Status of the valve on the top outlet during heating
\end{definition}
\begin{proof}
  This is the declared model definition of Valve Status.
\end{proof}
\begin{definition}[Figure Object]
  \label{def:physics:phyx-mini-0451:phyxminiproblems-problemphyxmini0451-figureobject}
  \lean{PhyXMiniProblems.ProblemPhyXMini0451.FigureObject}
  Distinct objects visible in the supplied raster
\end{definition}
\begin{proof}
  This is the declared model definition of Figure Object.
\end{proof}
\begin{definition}[Rigid Vessel Figure]
  \label{def:physics:phyx-mini-0451:phyxminiproblems-problemphyxmini0451-rigidvesselfigure}
  \lean{PhyXMiniProblems.ProblemPhyXMini0451.RigidVesselFigure}
  \uses{def:physics:phyx-mini-0451:phyxminiproblems-problemphyxmini0451-figureobject}
  Qualitative geometry transcribed from the primary image. The record contains no thermodynamic property value and no heat-transfer answer
\end{definition}
\begin{proof}
  This is the declared model definition of Rigid Vessel Figure.
\end{proof}
\begin{definition}[Water State]
  \label{def:physics:phyx-mini-0451:phyxminiproblems-problemphyxmini0451-waterstate}
  \lean{PhyXMiniProblems.ProblemPhyXMini0451.WaterState}
  \uses{def:physics:phyx-mini-0451:phyxminiproblems-problemphyxmini0451-specificvolumequantity, def:physics:phyx-mini-0451:phyxminiproblems-problemphyxmini0451-specificinternalenergyquantity, def:physics:phyx-mini-0451:phyxminiproblems-problemphyxmini0451-measuredtemperature, def:physics:phyx-mini-0451:phyxminiproblems-problemphyxmini0451-waterphaseregion}
  An equilibrium state of the fixed water inventory. vaporQuality is a dimensionless mass fraction whose thermodynamic meaning is imposed only when phaseRegion is the saturated liquid-vapor mixture
\end{definition}
\begin{proof}
  This is the declared model definition of Water State.
\end{proof}
\begin{definition}[Saturated Water Table]
  \label{def:physics:phyx-mini-0451:phyxminiproblems-problemphyxmini0451-saturatedwatertable}
  \lean{PhyXMiniProblems.ProblemPhyXMini0451.SaturatedWaterTable}
  \uses{def:physics:phyx-mini-0451:phyxminiproblems-problemphyxmini0451-specificvolumequantity, def:physics:phyx-mini-0451:phyxminiproblems-problemphyxmini0451-specificinternalenergyquantity, def:physics:phyx-mini-0451:phyxminiproblems-problemphyxmini0451-measuredtemperature}
  External saturated-water property data. The vaporization internal energy is u\_fg = u\_g - u\_f; all outputs remain dimensionful physical quantities
\end{definition}
\begin{proof}
  This is the declared model definition of Saturated Water Table.
\end{proof}
\begin{definition}[Rigid Water Vessel Setup]
  \label{def:physics:phyx-mini-0451:phyxminiproblems-problemphyxmini0451-rigidwatervesselsetup}
  \lean{PhyXMiniProblems.ProblemPhyXMini0451.RigidWaterVesselSetup}
  \uses{def:physics:phyx-mini-0451:phyxminiproblems-problemphyxmini0451-massquantity, def:physics:phyx-mini-0451:phyxminiproblems-problemphyxmini0451-volumequantity, def:physics:phyx-mini-0451:phyxminiproblems-problemphyxmini0451-workingfluid, def:physics:phyx-mini-0451:phyxminiproblems-problemphyxmini0451-processstage, def:physics:phyx-mini-0451:phyxminiproblems-problemphyxmini0451-processkind, def:physics:phyx-mini-0451:phyxminiproblems-problemphyxmini0451-valvestatus, def:physics:phyx-mini-0451:phyxminiproblems-problemphyxmini0451-rigidvesselfigure, def:physics:phyx-mini-0451:phyxminiproblems-problemphyxmini0451-waterstate, def:physics:phyx-mini-0451:phyxminiproblems-problemphyxmini0451-saturatedwatertable}
  The closed rigid-vessel experiment. Heat is positive into the water and boundary work is positive when done by the water. Endpoint observables are independent fields until related by the governing laws below
\end{definition}
\begin{proof}
  This is the declared model definition of Rigid Water Vessel Setup.
\end{proof}
\begin{definition}[Matches Rigid Water Vessel Scenario]
  \label{def:physics:phyx-mini-0451:phyxminiproblems-problemphyxmini0451-matchesrigidwatervesselscenario}
  \lean{PhyXMiniProblems.ProblemPhyXMini0451.MatchesRigidWaterVesselScenario}
  \uses{def:physics:phyx-mini-0451:phyxminiproblems-problemphyxmini0451-rigidwatervesselsetup}
  Qualitative physical idealizations supplied by the scenario
\end{definition}
\begin{proof}
  This is the declared model definition of Matches Rigid Water Vessel Scenario.
\end{proof}
\begin{definition}[Matches Supplied Figure]
  \label{def:physics:phyx-mini-0451:phyxminiproblems-problemphyxmini0451-matchessuppliedfigure}
  \lean{PhyXMiniProblems.ProblemPhyXMini0451.MatchesSuppliedFigure}
  \uses{def:physics:phyx-mini-0451:phyxminiproblems-problemphyxmini0451-figureobject, def:physics:phyx-mini-0451:phyxminiproblems-problemphyxmini0451-rigidwatervesselsetup}
  Facts visible in image 451.png: the vessel and inventory, four flames below it, and the top pipe with a crossed circular valve symbol. Closed-valve operation is kept in the scenario predicate rather than inferred solely from the symbol
\end{definition}
\begin{proof}
  This is the declared model definition of Matches Supplied Figure.
\end{proof}
\begin{definition}[Matches Problem Readouts]
  \label{def:physics:phyx-mini-0451:phyxminiproblems-problemphyxmini0451-matchesproblemreadouts}
  \lean{PhyXMiniProblems.ProblemPhyXMini0451.MatchesProblemReadouts}
  \uses{def:physics:phyx-mini-0451:phyxminiproblems-problemphyxmini0451-massinkilograms, def:physics:phyx-mini-0451:phyxminiproblems-problemphyxmini0451-temperatureindegreescelsius, def:physics:phyx-mini-0451:phyxminiproblems-problemphyxmini0451-rigidwatervesselsetup}
  Numerical measurements stated in the problem. The final temperature is recorded as a 20 °C rise rather than inserted directly as a table index. No heat value or answer label occurs here
\end{definition}
\begin{proof}
  This is the declared model definition of Matches Problem Readouts.
\end{proof}
\begin{definition}[Matches Reference Saturated Water Data]
  \label{def:physics:phyx-mini-0451:phyxminiproblems-problemphyxmini0451-matchesreferencesaturatedwaterdata}
  \lean{PhyXMiniProblems.ProblemPhyXMini0451.MatchesReferenceSaturatedWaterData}
  \uses{def:physics:phyx-mini-0451:phyxminiproblems-problemphyxmini0451-specificvolumeincubicmetresperkilogram, def:physics:phyx-mini-0451:phyxminiproblems-problemphyxmini0451-specificinternalenergyinkilojoulesperkilogram, def:physics:phyx-mini-0451:phyxminiproblems-problemphyxmini0451-rigidwatervesselsetup}
  Saturated-water table entries at the independently specified endpoint temperatures. Specific volumes are in m³/kg; specific internal energies are in kJ/kg. These reference data determine the final phase and quality but do not mention heat transfer or an answer choice
\end{definition}
\begin{proof}
  This is the declared model definition of Matches Reference Saturated Water Data.
\end{proof}
\begin{definition}[Satisfies Rigid Vessel Water Laws]
  \label{def:physics:phyx-mini-0451:phyxminiproblems-problemphyxmini0451-satisfiesrigidvesselwaterlaws}
  \lean{PhyXMiniProblems.ProblemPhyXMini0451.SatisfiesRigidVesselWaterLaws}
  \uses{def:physics:phyx-mini-0451:phyxminiproblems-problemphyxmini0451-massinkilograms, def:physics:phyx-mini-0451:phyxminiproblems-problemphyxmini0451-volumeincubicmetres, def:physics:phyx-mini-0451:phyxminiproblems-problemphyxmini0451-specificvolumeincubicmetresperkilogram, def:physics:phyx-mini-0451:phyxminiproblems-problemphyxmini0451-specificinternalenergyinkilojoulesperkilogram, def:physics:phyx-mini-0451:phyxminiproblems-problemphyxmini0451-energyinkilojoules, def:physics:phyx-mini-0451:phyxminiproblems-problemphyxmini0451-processstage, def:physics:phyx-mini-0451:phyxminiproblems-problemphyxmini0451-rigidwatervesselsetup}
  Governing laws for the closed rigid water system: mass is conserved and the occupied volume is unchanged; V = m v at each endpoint; a state whose specific volume lies between saturated liquid and vapor values at its temperature is a saturated mixture; mixture specific volume and internal energy interpolate by vapor quality; the closed-system first law is ΔU = Q - W\_by; a rigid boundary performs no boundary work. None of these laws includes the requested numerical heat or selects a multiple-choice answer
\end{definition}
\begin{proof}
  This is the declared model definition of Satisfies Rigid Vessel Water Laws.
\end{proof}
\begin{definition}[Answer Choice]
  \label{def:physics:phyx-mini-0451:phyxminiproblems-problemphyxmini0451-answerchoice}
  \lean{PhyXMiniProblems.ProblemPhyXMini0451.AnswerChoice}
  Labels printed beside the four candidate heat transfers
\end{definition}
\begin{proof}
  This is the declared model definition of Answer Choice.
\end{proof}
\begin{definition}[displayed Heat In Kilojoules]
  \label{def:physics:phyx-mini-0451:phyxminiproblems-problemphyxmini0451-displayedheatinkilojoules}
  \lean{PhyXMiniProblems.ProblemPhyXMini0451.displayedHeatInKilojoules}
  \uses{def:physics:phyx-mini-0451:phyxminiproblems-problemphyxmini0451-answerchoice}
  Heat-transfer value in kilojoules printed beside each answer label
\end{definition}
\begin{proof}
  This is the declared model definition of displayed Heat In Kilojoules.
\end{proof}
\begin{definition}[recorded Dataset Answer]
  \label{def:physics:phyx-mini-0451:phyxminiproblems-problemphyxmini0451-recordeddatasetanswer}
  \lean{PhyXMiniProblems.ProblemPhyXMini0451.recordedDatasetAnswer}
  \uses{def:physics:phyx-mini-0451:phyxminiproblems-problemphyxmini0451-answerchoice}
  Dataset answer metadata, deliberately not used as a theorem premise
\end{definition}
\begin{proof}
  This is the declared model definition of recorded Dataset Answer.
\end{proof}
\begin{definition}[Is Reported Heat Choice]
  \label{def:physics:phyx-mini-0451:phyxminiproblems-problemphyxmini0451-isreportedheatchoice}
  \lean{PhyXMiniProblems.ProblemPhyXMini0451.IsReportedHeatChoice}
  \uses{def:physics:phyx-mini-0451:phyxminiproblems-problemphyxmini0451-energyinkilojoules, def:physics:phyx-mini-0451:phyxminiproblems-problemphyxmini0451-rigidwatervesselsetup, def:physics:phyx-mini-0451:phyxminiproblems-problemphyxmini0451-answerchoice, def:physics:phyx-mini-0451:phyxminiproblems-problemphyxmini0451-displayedheatinkilojoules}
  A choice reports the modeled heat to the one-decimal precision of choice D. This remains a nontrivial relation between an independently modeled heat and the displayed table
\end{definition}
\begin{proof}
  This is the declared model definition of Is Reported Heat Choice.
\end{proof}
% --- Archon declaration topology end ---
% --- Archon physics formalization source end ---
